% !PRINTED 468
\weberpage{468}

%<<CONT
wird. Ist $\omega$ eine beliebige Grösse in $\Omega(\varrho)$, so kann man die mit $\omega$ conjugirten Grössen so darstellen:
\weq{4}{\omega = \varphi(\varrho),\ \omega_1 = \varphi(\varrho_1) \ldots \omega_{\mu-1} = \varphi(\varrho_{\mu-1}),}
worin $\varphi(t)$ eine rationale Function in $\Omega$ ist.

Da nun $g(t)$ irreducibel ist, so gilt der Satz \S.\,141, II, den wir jetzt so aussprechen:

1. \gesperrt{Wenn eine rationale Function $\Phi(t)$ in $\Omega$ eine Wurzel mit $g(t)$ gemeinsam hat, so verschwinden alle conjugirten Grössen}
\[
\Phi(\varrho),\ \Phi(\varrho_1) \ldots \Phi(\varrho_{\mu-1}).
\]

Wenn in einer der Functionen $\Theta_k(\varrho)$, durch die nach (3) die Wurzeln von (1) ausgedrückt sind, $\varrho$ durch eine andere Wurzel $\varrho_h$ ersetzt wird, so entsteht daraus wieder eine der Wurzeln; denn ist
\[
g[\Theta_k(\varrho)] = 0,
\]
so ist nach 1. auch
\[
g[\Theta_k(\varrho_h)] = 0,
\]
und die Reihe der Grössen
\weq{5}{\varrho_h,\ \Theta_1(\varrho_h),\ \Theta_2(\varrho_h) \ldots \Theta_{\mu-1}(\varrho_h)}
stimmt, von der Anordnung abgesehen, mit der Reihe
\weq{6}{\varrho,\ \Theta_1(\varrho),\ \Theta_2(\varrho) \ldots \Theta_{\mu-1}(\varrho)}
überein. Dies wird erwiesen sein, wenn wir zeigen, dass in (5) keine zwei gleichen Werthe vorkommen. Bezeichnen wir der Uebereinstimmung halber mit $\Theta_0(\varrho)$ oder $\Theta(\varrho)$ die Wurzel $\varrho$ selbst, so folgt aus der Gleichheit zweier der Grössen (5)
\weq{7}{\Theta_i(\varrho_h) = \Theta_k(\varrho_h)}
nach dem Satze 1.
\weq{8}{\Theta_i(\varrho) = \Theta_k(\varrho),}
was aber, wenn $i$ von $k$ verschieden ist, nicht möglich ist. Also sind zugleich mit den Grössen (6) auch die Grössen (5) unter einander verschieden.

Wir können dies als Satz zusammenfassen:

2. \gesperrt{Vertauscht man $\varrho$ mit $\varrho_h$, so geht zugleich jede mit $\varrho$ conjugirte Grösse in eine bestimmte andere über, und niemals zwei verschiedene in dieselbe.}

Wenn wir in allen Functionen von $\varrho$ statt $\varrho$ eine andere Wurzel $\varrho_a$ setzen, so führen wir eine \gesperrt{Substitution} aus. Wir bezeichnen diese Substitution mit
%CONT>>
